\begin{center}\begin{tikzpicture}[font=\small,>=stealth,line join=round]
\draw[gray!40,rounded corners=3pt] (0,0) rectangle (14.2,8.3);\node at (7.1,7.700000000000001) {算的是两个输入之差，不是中点到单侧的距离};
\draw[->,gray] (.8,3.2)--(13.1,3.2) node[right]{$x$};\foreach \v in {1,3,5,7,9}{\pgfmathsetmacro{\X}{1.8+\v}\draw[gray] (\X,3.2)--(\X,3.1);\node[below=7pt] at (\X,3.1){$\v$};}\foreach \v in {1,5,9}{\fill[blue] ({1.8+\v},3.2) circle(2.4pt);}\draw[<->,orange!85!black,thick] (2.8,4.7)--(6.8,4.7) node[midway,above]{$5-1=4$};\draw[<->,violet,thick] (6.8,6.1)--(10.8,6.1) node[midway,above]{$9-5=4$};\node at (7.1,1.8){$x=3$ 时，输入为 $x-2=1$ 与 $x+2=5$；中点到两边各2。};
\node[text width=12.6cm,align=center] at (7.1,.65) {相邻标点相隔4；仅凭这几个点，不能证明整个函数的周期。};
\end{tikzpicture}\end{center}
